\chapter{Operators and Expressions}%
\label{ch:operators}

So far our programs have mostly stored and printed values. Now they start to
\emph{compute}. An \textbf{operator} is a symbol that combines values to produce
a new one \code{+} adds, \code{<} compares, \code{\&\&} reasons about truth.
A combination of values and operators that produces a result is an
\textbf{expression}. This chapter is a tour of Salam's operators and the rules
that decide what an expression means.

\section{Arithmetic operators}

The five everyday arithmetic operators work as you would expect:

\begin{center}
\begin{tabular}{@{}cll@{}}
\toprule
\textbf{Operator} & \textbf{Meaning} & \textbf{Example} \\
\midrule
\code{+} & addition       & \code{3 + 4}  $\to$ \code{7} \\
\code{-} & subtraction    & \code{10 - 6} $\to$ \code{4} \\
\code{*} & multiplication & \code{6 * 7}  $\to$ \code{42} \\
\code{/} & division       & \code{17 / 5} $\to$ \code{3} \\
\code{\%} & remainder (modulo) & \code{17 \% 5} $\to$ \code{2} \\
\bottomrule
\end{tabular}
\end{center}

\begin{salam}
func main:
    println 3 + 4 * 2          // 11
    println 17 / 5, 17 % 5     // 3 2
    println 10 - 6, 6 * 7      // 4 42
end
\end{salam}

\begin{progout}
11
3 2
4 42
\end{progout}

\subsection{Integer versus floating-point division}

We met this in Chapter~2, but it is worth repeating because it trips up almost
everyone once. When \emph{both} operands of \code{/} are integers, you get
\emph{integer} division: the result is a whole number, with any fractional part
thrown away.

\begin{salam}
func main:
    println 7 / 2          // 3  (integer division)
    println 7.0 / 2.0      // 3.5 (floating-point division)
    println 7 as f64 / 2   // 3.5 (one float operand is enough)
end
\end{salam}

\begin{progout}
3
3.5
3.5
\end{progout}

The \code{\%} operator gives the \emph{remainder} of an integer division:
\code{7 \% 2} is \code{1}, because 7 is 2 threes with 1 left over. Modulo is
surprisingly useful: \code{n \% 2 == 0} tests whether \code{n} is even, and
\code{n \% 10} extracts the last decimal digit of \code{n}.

\subsection{The power operator}

Salam has a dedicated exponentiation operator, \code{**}. It raises the left
operand to the power of the right and produces a \emph{floating-point} result:

\begin{salam}
func main:
    println 2 ** 10        // 1024
    x : f64 = 9 ** 0.5       // 3 -- a square root, since x^(1/2) = sqrt(x)
    println x
end
\end{salam}

\begin{progout}
1024
3
\end{progout}

\begin{note}
Because \code{**} always yields a floating-point number, \code{2 ** 10} is the
floating value \code{1024.0}; it just prints as \code{1024} because there is no
fractional part to show. If you need an integer power and the values are small,
plain multiplication (or a loop, as in Chapter~5) keeps the result an integer.
\end{note}

\section{Compound assignment}

When you update a variable using its own current value a very common pattern ---
a \textbf{compound assignment} operator says it more briefly. Each one combines an
arithmetic operation with assignment:

\begin{center}
\begin{tabular}{@{}cl@{}}
\toprule
\textbf{Shorthand} & \textbf{Means} \\
\midrule
\code{x += y}  & \code{x = x + y} \\
\code{x -= y}  & \code{x = x - y} \\
\code{x *= y}  & \code{x = x * y} \\
\code{x /= y}  & \code{x = x / y} \\
\code{x \%= y} & \code{x = x \% y} \\
\code{x **= y} & \code{x = x ** y} \\
\bottomrule
\end{tabular}
\end{center}

\begin{salam}
func main:
    mut n : int = 10
    n += 5       // n is now 15
    n -= 2       // 13
    n *= 2       // 26
    println n

    mut p : f64 = 2.0
    p **= 10     // same as p = p ** 10
    println p
end
\end{salam}

\begin{progout}
26
1024
\end{progout}

\begin{warn}
Compound assignment only works on a \code{mut} variable it \emph{changes} the
variable, after all. Used on an immutable binding, it is the same immutability
error you saw in Chapter~2.
\end{warn}

Because \code{**} always produces a floating-point result (see the previous
section), \code{x **= y} only type-checks when \code{x} itself is a \code{f64}
(or a \code{f32}, or a type with an \code{operator **} overload from
Chapter~13) for the same reason \code{x = x ** y} would need \code{x} to accept
a float on the right-hand side.

\begin{note}
Every compound assignment is checked with exactly the same rules as writing the
operation out longhand. \code{x \%= y} still requires both sides to be integers,
and \code{x += y} still refuses to mix a \code{bool} with anything numeric, even
though \code{x} and \code{y} are the same type. In other words, \code{x op= y}
is never more permissive than \code{x = x op y}; it is purely a shorter way to
write it.
\end{note}

\section{Comparison operators}

A comparison asks a yes/no question about two values and answers with a
\code{bool}. There are six:

\begin{center}
\begin{tabular}{@{}cll@{}}
\toprule
\textbf{Operator} & \textbf{Question} & \textbf{Example} \\
\midrule
\code{==} & equal?                 & \code{5 == 5} $\to$ \code{true} \\
\code{!=} & not equal?             & \code{5 != 3} $\to$ \code{true} \\
\code{<}  & less than?             & \code{3 < 5}  $\to$ \code{true} \\
\code{>}  & greater than?          & \code{3 > 5}  $\to$ \code{false} \\
\code{<=} & less than or equal?    & \code{5 <= 5} $\to$ \code{true} \\
\code{>=} & greater than or equal? & \code{4 >= 5} $\to$ \code{false} \\
\bottomrule
\end{tabular}
\end{center}

\begin{salam}
func main:
    println 5 > 3, 5 == 5, 5 != 3, 5 <= 5
end
\end{salam}

\begin{progout}
true true true true
\end{progout}

\begin{warn}
Do not confuse \code{=} with \code{==}. A single \code{=} \emph{assigns} a value
to a variable; a double \code{==} \emph{compares} two values and yields a
\code{bool}. Writing \code{=} where you meant \code{==} is a classic slip ---
fortunately Salam's static checks catch most such mistakes for you.
\end{warn}

Comparisons also work on strings: \code{==} and \code{!=} test whether two
strings have the same contents (not merely whether they are the same object),
and \code{<}, \code{>}, and friends compare them in dictionary order.

\begin{salam}
func main:
    println "abc" == "abc"     // true
    println "abc" == "abd"     // false
end
\end{salam}

\begin{progout}
true
false
\end{progout}

\section{Logical operators}

Logical operators combine \code{bool} values the raw material of decisions.
There are three:

\begin{center}
\begin{tabular}{@{}cll@{}}
\toprule
\textbf{Operator} & \textbf{Meaning} & \textbf{True when\dots} \\
\midrule
\code{\&\&} & and & both sides are true \\
\code{||}   & or  & at least one side is true \\
\code{!}    & not & the operand is false \\
\bottomrule
\end{tabular}
\end{center}

\begin{salam}
func main:
    age : int = 20
    has_ticket : bool = true
    can_enter : auto = age >= 18 && has_ticket
    println can_enter
    println true || false, !true
end
\end{salam}

\begin{progout}
true
true false
\end{progout}

\begin{deep}[In Depth: short-circuit evaluation]
\code{\&\&} and \code{||} are \emph{short-circuiting}: they evaluate the
right-hand side only if they must. In \code{a \&\& b}, if \code{a} is already
\code{false}, the whole expression is \code{false} regardless of \code{b}, so
\code{b} is never evaluated. Likewise \code{a || b} skips \code{b} when \code{a}
is \code{true}. This is not just an optimization it lets you write a guard and
a use of it in one line, e.g. testing that a count is positive \emph{before}
dividing by it, knowing the risky part runs only when it is safe.
\end{deep}

\begin{note}
Salam spells the logical operators with symbols (\code{\&\&}, \code{||},
\code{!}), not words. There is no \code{and}, \code{or}, or \code{not} keyword.
\end{note}

\section{A note on bitwise operators}

Some languages provide \emph{bitwise} operators (\code{\&}, \code{|},
\code{<<}, and so on) that manipulate the individual bits of an integer. The
current version of Salam does \emph{not} include them writing \code{6 | 1},
for instance, is a syntax error (the compiler even reminds you that \code{||} is
the logical ``or''). If you need bit manipulation today, you can fall back to
arithmetic (multiplying or dividing by powers of two shifts bits) or call a C
function through the foreign function interface of Chapter~19.

\section{String concatenation}

You have already seen that \code{+} joins two strings end to end. This is the
same \code{+} symbol as numeric addition, but Salam picks the right meaning from
the types of the operands: two numbers add, two strings concatenate.

\begin{salam}
func main:
    first : str = "Salam"
    full : auto = first + ", " + "world!"
    println full
end
\end{salam}

\begin{progout}
Salam, world!
\end{progout}

We devote all of Chapter~8 to strings; for now, just remember that \code{+}
between two \code{str} values builds a new, larger string.

\section{Precedence and grouping}

When several operators appear in one expression, \textbf{precedence} decides who
acts first. The rules follow ordinary mathematics: \code{**} binds tightest, then
\code{*}, \code{/}, and \code{\%}, then \code{+} and \code{-}, then the
comparisons, and finally the logical operators. So

\begin{salam}
func main:
    println 3 + 4 * 2          // 11, not 14: * happens before +
    println 2 + 3 > 4          // true: 2+3 first, then compare
end
\end{salam}

\begin{progout}
11
true
\end{progout}

When in doubt or simply to make your intent unmistakable add parentheses.
They cost nothing and override any precedence rule:

\begin{salam}
func main:
    println (3 + 4) * 2        // 14: parentheses force + first
end
\end{salam}

\begin{progout}
14
\end{progout}

\begin{tip}
You do not need to memorise the full precedence table. Two habits cover almost
everything: remember that multiplication and division come before addition and
subtraction, and reach for parentheses whenever an expression is even slightly
hard to read at a glance. Clear beats clever.
\end{tip}

\section{Putting it together}

Here is a small program that uses several operator families at once: it decides
whether a year is a leap year, a genuinely useful rule built entirely from
\code{\%}, comparison, and logical operators.

\begin{salam}
func main:
    year : int = 2024
    leap : auto = (year % 4 == 0 && year % 100 != 0) || (year % 400 == 0)
    println year, "leap?", leap
end
\end{salam}

\begin{progout}
2024 leap? true
\end{progout}

Read the condition in English: a year is a leap year if it is divisible by 4 but
\emph{not} by 100 unless it is also divisible by 400. The operators let us say
that precisely. In Chapter~4 we will use such conditions to make the program
actually \emph{do} different things.

\section{Chapter summary}

\begin{itemize}[leftmargin=1.4em]
  \item Arithmetic: \code{+ - * / \%}, plus \code{**} for powers (which yields a
        floating-point result).
  \item Integer \code{/} truncates; use a float operand for fractional results.
        \code{\%} gives the remainder.
  \item Compound assignment (\code{+=}, \code{-=}, \code{*=}, \code{/=},
        \code{\%=}, \code{**=}) updates a \code{mut} variable in place, and is
        checked with exactly the same type rules as its longhand form.
  \item Comparisons (\code{== != < > <= >=}) yield a \code{bool} and also work on
        strings.
  \item Logical operators are \code{\&\&}, \code{||}, \code{!}, and they
        short-circuit. There are no word forms and no bitwise operators.
  \item \code{+} concatenates strings.
  \item Precedence follows mathematics; parentheses override it and aid clarity.
\end{itemize}

\section{Exercises}

\exercise Print the results of \code{20 / 3}, \code{20 \% 3}, and
\code{20.0 / 3.0}. Explain in a comment why the first two differ from the third.

\exercise Given \code{mut x int = 100}, use compound assignment to halve it, then
add 7, then print it.

\exercise Write a boolean expression that is \code{true} exactly when an integer
\code{n} is between 1 and 12 inclusive, and test it with a couple of values.

\exercise Without running it first, predict the value of
\code{2 + 3 * 4 > 10 \&\& 1 < 2}. Then run it and check.

\exercise Build a person's full name from \code{first} and \code{last} string
variables with a space between them, using \code{+}, and print it.

\exercise Given \code{mut side : f64 = 3.0}, use \code{**=} to compute the cube
of \code{side} in place, then print it. What error do you get if you try the
same thing starting from \code{mut side : int = 3} instead? Why?
